\documentclass{article} % For LaTeX2e
\usepackage{iclr2026_conference,times}

\usepackage{amsmath,amssymb,amsfonts,amsthm,mathtools}
\usepackage{booktabs}
\usepackage{graphicx}
\usepackage{enumitem}
\usepackage{xcolor}
\usepackage{hyperref}
\usepackage{url}
\usepackage{microtype}
\usepackage{wrapfig}

\hypersetup{colorlinks=true,linkcolor=blue,citecolor=blue,urlcolor=blue}

\newtheorem{theorem}{Theorem}
\newtheorem{proposition}{Proposition}
\newtheorem{lemma}{Lemma}
\newtheorem{corollary}{Corollary}
\newtheorem{assumption}{Assumption}
\newtheorem{definition}{Definition}
\newtheorem{remark}{Remark}

\newcommand{\E}{\mathbb{E}}
\newcommand{\R}{\mathbb{R}}
\newcommand{\diag}{\operatorname{diag}}
\newcommand{\tr}{\operatorname{tr}}
\newcommand{\rank}{\operatorname{rank}}
\newcommand{\as}{\asymp}
\newcommand{\pre}{\preceq}
\newcommand{\suc}{\succeq}
\newcommand{\lam}{\lambda}
\newcommand{\Sig}{\Sigma}
\newcommand{\tSig}{\widetilde\Sigma}
\newcommand{\tlam}{\widetilde\lambda}
\newcommand{\tmu}{\widetilde\mu}
\newcommand{\ts}{\widetilde s}
\newcommand{\calF}{\mathcal F}
\newcommand{\calE}{\mathcal E}
\newcommand{\calD}{\mathcal D}
\newcommand{\eff}{\mathrm{eff}}
\newcommand{\qeff}{q_{\mathrm{eff}}}
\newcommand{\Neff}{N_{\mathrm{eff}}}
\newcommand{\Kvis}{K_{\rho,\theta}}
\newcommand{\wd}{\mathrm{wd}}
\DeclareMathOperator*{\argmin}{arg\,min}

\title{Visible Spectrum: When Adaptive Optimizers Change Scaling Laws}

% Authors are hidden in submission mode. Uncomment \iclrfinalcopy only for the
% camera-ready version.
\author{Jingxu Xie et al.}

%\iclrfinalcopy
\begin{document}
\maketitle

\begin{abstract}
Modern scaling laws forecast how loss improves with data, parameters, and compute, but usually treat the optimizer as a constant-factor detail.  This paper shows that this misses a fundamental effect: even in Gaussian linear regression, coordinatewise adaptive optimizers change scaling exponents only when their second moments make the data spectrum visible.  We introduce a \emph{visible spectrum} exponent $\theta\in[0,1]$.  If the diagonal preconditioner learned by RMSProp/Adam is spectrally comparable to $(\Sigma^\theta+\rho I)^{-1/2}$, then the preconditioned covariance satisfies
\[
    \lambda_i(P^{1/2}\Sigma P^{1/2})
    \asymp
    \lambda_i(\Sigma)(\lambda_i(\Sigma)^\theta+\rho)^{-1/2},
    \qquad
    \qeff=\theta/2.
\]
Thus, for power-law spectra $\lambda_i\asymp i^{-a}$, the learned-mode count changes from the SGD rate $n^{1/a}$ to $n^{1/[a(1-\theta/2)]}$ before damping.  With source exponent $b$, this yields sharp risk and compute laws, including the compute exponent
\[
    \beta(\theta)=\frac{b-1}{\max\{a(1-\theta/2),b\}+1}.
\]
The theory separates the roles of Adam's components: second moments determine the spectral exponent, first-moment momentum changes constants and stability, and AdamW weight decay imposes a learned-mode cap that must decay with compute to preserve the no-decay law.  Experiments with deterministic filters, stochastic training, coordinate rotations, random-feature/sketched covariances, compute-optimal sweeps, and learning-rate grids support the mechanism.  The results suggest that adaptive optimizers change scaling laws through \emph{visible spectral information}, not through adaptivity alone.
\end{abstract}

\section{Introduction}

Scaling laws have become a central tool for designing modern machine learning systems.  They inform model-size and data-allocation decisions, motivate compute-optimal training recipes, and provide empirical regularities that often hold across orders of magnitude \citep{hestness2017deep,kaplan2020scaling,henighan2020scaling,hoffmann2022training}.  Yet scaling laws are usually stated as properties of model families, datasets, and compute budgets, with the optimizer treated as a tuning detail.  This is unsatisfactory: most frontier models are trained with Adam or AdamW rather than vanilla SGD, and practical performance is often highly optimizer-dependent.

This paper asks a deliberately narrow but fundamental question:
\begin{quote}
    \emph{When can an adaptive optimizer change a scaling exponent, rather than merely a constant?}
\end{quote}
We answer this question in power-law linear regression.  This setting is simple enough to prove sharp statements, but rich enough to exhibit the same ingredients that appear in modern scaling-law theory: approximation error, optimization bias, variance/effective dimension, power-law spectra, source conditions, compute allocation, and optimizer-induced preconditioning \citep{bartlett2020benign,hastie2022surprises,bordelon2020spectrum,canatar2021spectral,lin2024scaling}.

\paragraph{The short answer: adaptivity helps exponents only through visible spectrum.}
A coordinatewise adaptive method does not directly observe covariance eigenvalues.  It observes coordinatewise gradient second moments.  For Gaussian squared loss, these second moments are proportional to the diagonal entries of the covariance in the optimizer coordinate system.  Therefore RMSProp/Adam acts spectrally only if those diagonal entries retain spectral information.  If optimizer coordinates are eigenvectors, the second moment in coordinate $i$ is proportional to $\lambda_i$, and Adam/RMSProp behaves like the damped spectral preconditioner $(\Sigma+\rho I)^{-1/2}$.  If the coordinates are globally mixed so that all coordinate variances are nearly equal, the same coordinatewise rule is scalar up to constants and cannot change the power-law spectral exponent.

\paragraph{Visible spectrum.}
We formalize this by a visible spectral exponent $\theta\in[0,1]$.  Let $z\sim\mathcal N(0,\Sigma)$ be the trainable feature vector and let
\[
    P=\diag((\Sigma_{jj}+\rho)^{-1/2})
\]
be the idealized diagonal RMSProp/Adam preconditioner, ignoring scalar residual factors.  If
\begin{equation}
\label{eq:intro-visible}
    P\asymp (\Sigma^\theta+\rho I)^{-1/2}
\end{equation}
in Loewner order, then the preconditioned covariance has eigenvalues
\[
    \lambda_i(P^{1/2}\Sigma P^{1/2})
    \asymp
    \lambda_i(\Sigma)(\lambda_i(\Sigma)^\theta+\rho)^{-1/2}.
\]
For $\lambda_i(\Sigma)\asymp i^{-a}$, this changes the effective exponent from $a$ to $a(1-\theta/2)$ before the damping knee.  We call $\qeff=\theta/2$ the optimizer's effective spectral preconditioning exponent.

\paragraph{Why this matters.}
This result has two consequences that are easy to conflate.  First, when $\theta>0$, coordinatewise adaptivity can change the asymptotic learned-mode count and the compute-optimal exponent.  Second, even when $\theta=0$, adaptive optimizers may still improve finite-time risk through scalar normalization, momentum, or stability.  Our experiments show both phenomena.  In globally mixed Haar and Gaussian-sketch coordinates, Adam/RMSProp can beat SGD in finite time, but their measured $\qeff\approx0$; a spectral oracle with $\qeff=1/2$ then wins after a wider learning-rate sweep.  Thus the right question is not whether Adam ``helps,'' but whether its second moments expose the spectrum well enough to change the exponent class.

\paragraph{Contributions.}
\begin{enumerate}[leftmargin=1.4em,itemsep=0.25em]
    \item We prove a \textbf{visible-spectrum theorem}: spectral comparability \eqref{eq:intro-visible} implies $\qeff=\theta/2$ and determines the preconditioned spectrum.
    \item We derive sharp risk laws and compute-optimal allocations for visible-spectrum preconditioning, including a visibility phase transition where increasing $\theta$ no longer improves the compute exponent.
    \item We decompose Adam/RMSProp/AdamW: the second moment determines $\qeff$, first-moment momentum changes constants and stability, and AdamW weight decay adds a learned-mode cap requiring a compute-dependent schedule to preserve scaling.
    \item We prove alignment results: eigenbasis and band-limited features expose spectrum; flat, Haar, and isotropic Gaussian sketches hide it.  We also prove transformed-source stability for spectral and invariant-band preconditioners.
    \item We validate the theory through deterministic filter sweeps, stochastic training, learning-rate grids, visible-profile interpolation, source diagnostics, compute-optimal sweeps, AdamW decay schedules, and random-feature/sketched-coordinate experiments.
\end{enumerate}

\begin{figure}[t]
\centering
\fbox{\begin{minipage}{0.93\linewidth}
\textbf{Placeholder for Figure 1: the visible-spectrum mechanism.}  Left: coordinatewise second moments track $\diag(\Sigma)$, not eigenvalues directly.  Middle: aligned or band-limited features make $P\approx(\Sigma+\rho I)^{-1/2}$ and yield $\qeff\approx1/2$; flat/Haar/global sketches make $P\approx cI$ and yield $\qeff\approx0$.  Right: the visible exponent $\theta$ controls learned-mode count, risk, compute allocation, and AdamW decay schedules.
\end{minipage}}
\caption{Conceptual overview.  Coordinatewise adaptivity changes scaling exponents only when the coordinate basis makes spectral information visible.}
\label{fig:overview}
\end{figure}

\section{Setup}

Let $x\sim\mathcal N(0,H)$ and $y=\langle x,w^\star\rangle+\xi$, with $\E\xi=0$ and $\E\xi^2=\sigma^2$.  We study a trainable feature representation $z=Sx$ with covariance
\[
    \Sigma=\E[zz^\top]=SHS^\top.
\]
This includes the diagonal eigenbasis model, orthogonal coordinate changes, and sketched/random-feature representations.  We measure excess prediction risk in the trainable representation:
\[
    R(w)-\sigma^2 = \|w-w^\star\|_\Sigma^2.
\]
Let $\lambda_i(\Sigma)$ denote the ordered eigenvalues of $\Sigma$.

\begin{assumption}[Power-law spectrum]
\label{ass:spectrum}
For some $a>1$,
\[
    \lambda_i(\Sigma)\asymp i^{-a}.
\]
\end{assumption}

Let $\Sigma v_i=\lambda_i v_i$ and define source energies
\[
    s_i=\lambda_i\langle v_i,w^\star\rangle^2.
\]

\begin{assumption}[Source condition]
\label{ass:source}
For some $b>1$,
\[
    s_i\asymp i^{-b}.
\]
For noncommuting preconditioners, our risk theorem assumes the analogous condition in the transformed basis.  Appendix~\ref{app:source} shows that this is automatic for exact spectral preconditioners and holds in band-averaged form for invariant band-limited decompositions.
\end{assumption}

We write $N$ for the number of samples/updates, $M$ for the model dimension, and $\Neff=N/\log N$ for the usual effective one-pass horizon.  The scalar effective optimization horizon is denoted $n=\Neff\gamma$, with scalar normalizations from residual energy absorbed into $\gamma$ and $\rho$.

\section{The visible spectrum of coordinatewise adaptivity}

\subsection{Gradient second moments reveal coordinate variances}

\begin{lemma}[Coordinate second moments]
\label{lem:grad-second}
Let $z\sim\mathcal N(0,\Sigma)$ and $y=\langle z,w^\star\rangle+\xi$.  For error $e=w-w^\star$ and squared-loss gradient
\[
    g=(\langle z,e\rangle-\xi)z,
\]
we have, for every coordinate $j$,
\[
    \E[g_j^2\mid e]
    =
    c_e\Sigma_{jj}+2(\Sigma e)_j^2,
    \qquad
    c_e=\sigma^2+\|e\|_\Sigma^2.
\]
Consequently,
\[
    c_e\Sigma_{jj}
    \le
    \E[g_j^2\mid e]
    \le
    3c_e\Sigma_{jj}.
\]
\end{lemma}

\begin{proof}
Let $U=\langle z,e\rangle-\xi$ and $Z_j=z_j$.  Conditionally on $e$, $(U,Z_j)$ is centered Gaussian with
\[
    \E U^2=c_e,
    \qquad
    \E Z_j^2=\Sigma_{jj},
    \qquad
    \E[UZ_j]=(\Sigma e)_j.
\]
Isserlis' identity gives $\E[U^2Z_j^2]=\E U^2\E Z_j^2+2\E[UZ_j]^2$, proving the identity.  The upper bound follows from Cauchy-Schwarz in the $\Sigma$-inner product:
\[
    (\Sigma e)_j^2\le(e^\top\Sigma e)\Sigma_{jj}\le c_e\Sigma_{jj}.
\]
\end{proof}

Thus an RMSProp/Adam second-moment estimate tracks $\Sigma_{jj}$ up to a scalar residual-energy factor.  That scalar changes the learning rate and damping, but not the spectral exponent.  This observation is the bridge between optimizer dynamics and feature-coordinate geometry.

\subsection{Visible spectral comparability}

\begin{definition}[Visible spectral exponent]
\label{def:visible}
A diagonal adaptive preconditioner $P\succ0$ has visible exponent $\theta\in[0,1]$ at damping $\rho\ge0$ if there are constants $0<c\le C<\infty$, independent of dimension and scale, such that
\begin{equation}
\label{eq:comparability}
    c(\Sigma^\theta+
ho I)^{-1/2}
    \pre
    P
    \pre
    C(\Sigma^\theta+
ho I)^{-1/2}.
\end{equation}
\end{definition}

The case $\theta=1$ is true spectral Adam/RMSProp preconditioning.  The case $\theta=0$ is scalar preconditioning.  Intermediate $\theta$ means that coordinatewise second moments expose a fractional spectral profile.

\begin{theorem}[Visible-spectrum theorem]
\label{thm:visible-spectrum}
Assume \eqref{eq:comparability}.  Let $\tSig=P^{1/2}\Sigma P^{1/2}$.  Then
\[
    \lambda_i(\tSig)
    \asymp
    \lambda_i(\Sigma)\bigl(\lambda_i(\Sigma)^\theta+
ho\bigr)^{-1/2}.
\]
In particular, if $\lambda_i(\Sigma)\asymp i^{-a}$ and damping is inactive on the range of interest, then
\[
    \lambda_i(\tSig)\asymp i^{-a(1-	heta/2)},
    \qquad
    \qeff=\theta/2.
\]
\end{theorem}

\begin{proof}
The nonzero eigenvalues of $P^{1/2}\Sigma P^{1/2}$ equal those of $\Sigma^{1/2}P\Sigma^{1/2}$.  Applying \eqref{eq:comparability},
\[
    c\,\Sigma^{1/2}(\Sigma^\theta+
ho I)^{-1/2}\Sigma^{1/2}
    \pre
    \Sigma^{1/2}P\Sigma^{1/2}
    \pre
    C\,\Sigma^{1/2}(\Sigma^\theta+
ho I)^{-1/2}\Sigma^{1/2}.
\]
The outer terms commute with $\Sigma$ and have eigenvalues $\lambda_i(\Sigma)(\lambda_i(\Sigma)^\theta+
ho)^{-1/2}$.  The min-max principle gives the result.
\end{proof}

\begin{corollary}[Learned-mode count]
\label{cor:learned-count}
Let
\[
    K_{\rho,\theta}(n)=\#\{i:\lambda_i(\tSig)\gtrsim n^{-1}\}.
\]
For $\theta>0$,
\[
    K_{\rho,\theta}(n)
    \asymp
    \begin{cases}
    n^{1/[a(1-	heta/2)]}, & n\lesssim \rho^{-(1/\theta-1/2)},\\[4pt]
    \rho^{-1/(2a)}n^{1/a}, & n\gtrsim \rho^{-(1/\theta-1/2)}.
    \end{cases}
\]
For $\theta=0$, $K_{\rho,0}(n)\asymp n^{1/a}$.
\end{corollary}

\paragraph{Interpretation.}
The theorem converts a property of optimizer coordinates into a scaling-law exponent.  Aligned and band-limited coordinates expose $\theta\approx1$ and hence $\qeff\approx1/2$.  Flat, Haar-mixed, or isotropic Gaussian-sketch coordinates often expose $\theta\approx0$ and hence $\qeff\approx0$.  This is not a statement about whether Adam improves finite-time risk; it is a statement about whether coordinatewise adaptivity changes the spectral power law.

\section{Risk, compute, and optimizer components}

\subsection{Risk law}

The visible-spectrum theorem reduces the analysis to one-pass SGD on a transformed covariance.  Let $\tmu_i=\lambda_i(\tSig)$ and let $\ts_i$ denote the source energy in the transformed basis.  The exact filter form is
\begin{equation}
\label{eq:filters}
    \E R_M(w_N)-\sigma^2
    \asymp
    \sum_{i>M}\ts_i
    +
    \sum_{i\le M}\frac{\ts_i}{1+n\tmu_i}
    +
    \frac1{\Neff}\sum_{i\le M}\min\{1,n^2\tmu_i^2\}.
\end{equation}
This is the form used in the proofs.  Under the power-law and source assumptions it simplifies to a learned-count law.

\begin{theorem}[Visible-spectrum risk law]
\label{thm:risk}
Let $\alpha=a(1-	heta/2)$.  Assume the transformed source condition $\ts_i\asymp i^{-b}$ and the clean range $1<b<\alpha+1$.  Then, up to constants and logarithmic factors,
\[
    \E R_M(w_N)-\sigma^2
    \asymp
    M^{1-b}
    +K_{\rho,\theta}(n)^{1-b}
    +\frac{\min\{M,K_{\rho,\theta}(n)\}}{\Neff}.
\]
\end{theorem}

\begin{remark}[Why the source condition is not cosmetic]
In noncommuting coordinate systems, the transformed eigenvectors can rotate the target.  Appendix~\ref{app:source} proves source stability for exact spectral preconditioners and invariant band decompositions, and our source diagnostics test the condition in globally mixed cases.
\end{remark}

\subsection{Compute-optimal scaling}

Let compute be $C\asymp MN$.  In the pre-damping regime, write $K$ for the number of learned modes.  The sharp proxy is
\begin{equation}
\label{eq:compute-proxy}
    R(M,N,K)-\sigma^2
    \asymp
    M^{1-b}+K^{1-b}+K/N,
\end{equation}
subject to the algorithmic constraint $K\le N^{1/\alpha}$, where $\alpha=a(1-	heta/2)$.

\begin{theorem}[Compute-optimal visible scaling]
\label{thm:compute}
Let
\[
    m=\max\{\alpha,b\},
    \qquad
    \alpha=a(1-	heta/2).
\]
Then
\[
    M_\star(C)\asymp C^{1/(m+1)},
    \qquad
    N_\star(C)\asymp C^{m/(m+1)},
    \qquad
    K_\star(C)\asymp C^{1/(m+1)}.
\]
Moreover,
\[
    R_\star(C)-\sigma^2
    \asymp
    C^{-(b-1)/(m+1)},
\]
so the compute-risk exponent is
\[
    \beta(\theta)=\frac{b-1}{\max\{a(1-	heta/2),b\}+1}.
\]
\end{theorem}

\begin{corollary}[Visibility phase transition]
\label{cor:phase}
Define $\theta_c=2(1-b/a)$.  When $0<\theta_c<1$,
\[
    \beta(\theta)=
    \begin{cases}
    \dfrac{b-1}{a(1-	heta/2)+1}, & 0\le\theta\le\theta_c,\\[8pt]
    \dfrac{b-1}{b+1}, & \theta\ge\theta_c.
    \end{cases}
\]
Thus increasing visible spectral information improves the compute exponent only until the variance-limited phase, after which it improves constants and the optimal learning-rate schedule but not the asymptotic compute exponent.
\end{corollary}

\subsection{Adam, momentum, and AdamW}

The theory also separates the pieces of Adam.

\paragraph{Second moments.}
Lemma~\ref{lem:grad-second} implies that RMSProp/Adam second moments learn coordinate variances.  This determines $P$ and therefore $\qeff$.

\paragraph{First-moment momentum.}
For a fixed preconditioner with effective eigenvalue $\mu_i$, Adam-style first-moment momentum has population error polynomial
\[
    p_{t+1}^{(\beta_1)}(x)
    =
    (1+\beta_1-(1-\beta_1)x)p_t^{(\beta_1)}(x)
    -\beta_1p_{t-1}^{(\beta_1)}(x),
    \qquad x=\gamma\mu_i.
\]
The root near one is $1-x+O_{\beta_1}(x^2)$, so the cutoff remains $N\gamma\mu_i\asymp1$.  Momentum changes constants, transient behavior, and stability, but not the spectral power.

\paragraph{AdamW weight decay.}
Decoupled weight decay adds
\[
    w_{t+1}=(1-\gamma\delta)w_t-\gamma P_t m_t.
\]
For fixed effective eigenvalue $\mu_i$, the population error has a shrinkage floor $\delta/(\mu_i+\delta)$.  Hence learned modes must satisfy $\mu_i\gtrsim\max\{n^{-1},\delta\}$.  If $\delta(C)=C^{-s}$, then $K\le C^{s/\alpha}$.

\begin{theorem}[Compute-dependent AdamW decay]
\label{thm:adamw}
Let $m=\max\{\alpha,b\}$ and $\delta(C)=C^{-s}$.  The compute-risk exponent under the weight-decay cap is
\[
    \beta_{\rm wd}(\theta,s)
    =
    \min\left\{
        \frac{(b-1)s}{\alpha},
        \frac{b-1}{m+1}
    \right\}.
\]
The no-decay compute law is recovered iff
\[
    s\ge \frac{\alpha}{m+1}.
\]
Fixed weight decay corresponds to $s=0$ and creates an asymptotic floor.
\end{theorem}

\section{When is the spectrum visible?}

The visible-spectrum theorem is useful only if we can identify when \eqref{eq:comparability} holds.  We give three cases.

\begin{proposition}[Aligned coordinates]
If optimizer coordinates are eigenvectors of $\Sigma$, then $P=(\Sigma+
ho I)^{-1/2}$ and $\theta=1$.
\end{proposition}

\begin{proposition}[Band-limited mixing]
Suppose coordinates are partitioned into spectral bands $B_\ell$ such that $\lambda_i\asymp\Lambda_\ell$ for $i\in B_\ell$, and the optimizer basis mixes coordinates only within each band.  Then
\[
    P\asymp (\Sigma+
ho I)^{-1/2}
\]
up to band constants, so $\theta=1$ and $\qeff=1/2$.
\end{proposition}

\begin{proposition}[Flat and globally mixed coordinates]
If the coordinate variances satisfy $d_{\min}\le\Sigma_{jj}\le d_{\max}$ with $d_{\max}/d_{\min}=O(1)$, then $P\asymp cI$ and the spectrum of $P^{1/2}\Sigma P^{1/2}$ is the spectrum of $\Sigma$ up to constants.  Thus $\theta=0$ and $\qeff=0$.
\end{proposition}

For a Haar-random global rotation or an isotropic Gaussian sketch with sufficiently concentrated coordinate variances, the flat-coordinate condition holds approximately.  Conversely, structured or band-limited random features can preserve spectral visibility.

\begin{figure}[t]
\centering
\fbox{\begin{minipage}{0.93\linewidth}
\textbf{Placeholder for Figure 2: visibility regimes.} Plot measured $\qeff$ across feature systems: aligned, band-limited, synthetic $\theta$ profiles, Haar, and Gaussian sketch.  Expected: aligned/band near $0.5$, synthetic profiles on the line $\theta/2$, Haar/sketch near $0$.
\end{minipage}}
\caption{Feature coordinates determine whether coordinatewise second moments are spectral or scalar.}
\label{fig:qeff-regimes}
\end{figure}

\section{Experiments}

Our experiments are synthetic by design.  Each isolates one mechanism in the theory.  The recurring diagnostics are: (i) fitted $\qeff$, obtained by comparing the power-law slopes of the original and preconditioned spectra; (ii) excess-risk curves under actual mini-batch training; and (iii) compute-optimal exponents from the sharp risk proxy.

\subsection{Deterministic filters and visibility diagnostics}
Deterministic filter sweeps verify the learned-mode predictions for SGD, Adam/RMSProp-like spectra, and AdamW weight decay.  Coordinate-visibility sweeps show that aligned coordinates have effective spectral slope $-a/2$, while flat/Haar coordinates retain slope $-a$.  These runs validate the spectral algebra before adding stochastic optimizer noise.

\subsection{Stochastic training and tuned comparisons}
Mini-batch training confirms that RMSProp/Adam/AdamW second moments learn $\qeff\approx1/2$ in aligned and band-limited coordinates, while flat/Haar coordinates have $\qeff\approx0$.  Learning-rate sweeps show that flat/Haar Adam can improve finite-time risk despite $\qeff\approx0$.  We therefore report both best final risk and $\qeff$: risk captures finite-time performance, while $\qeff$ identifies the scaling-law mechanism.

\subsection{Visible-profile interpolation}
A controlled stochastic experiment uses fixed preconditioners
\[
    P_{\theta,i}=(\lambda_i^\theta+
ho)^{-1/2}.
\]
For $a=1.5,b=1.4$, tuned final risk decreases monotonically as $\theta$ increases, and the fitted risk slopes steepen accordingly (Table~\ref{tab:theta}).

\begin{table}[t]
\centering
\caption{Stochastic visible-profile training.  Increasing visible exponent $\theta$ improves tuned final risk and steepens the fitted risk slope.}
\label{tab:theta}
\begin{tabular}{ccccc}
\toprule
$\theta$ & best LR & final risk & risk slope & predicted $\qeff$ \\
\midrule
0.00 & 0.03 & 0.3177 & $-0.330$ & 0.000 \\
0.25 & 0.03 & 0.2245 & $-0.399$ & 0.125 \\
0.50 & 0.03 & 0.1330 & $-0.509$ & 0.250 \\
0.75 & 0.03 & 0.0534 & $-0.724$ & 0.375 \\
1.00 & 0.03 & 0.0074 & $-1.252$ & 0.500 \\
\bottomrule
\end{tabular}
\end{table}

\subsection{Compute-optimal scaling and AdamW schedules}
Compute-optimal sweeps optimize \eqref{eq:compute-proxy} under $C=MN$.  In a saturating case ($a=1.5,b=1.4$), the risk exponent plateaus after crossing the visibility phase boundary.  In a hard-source full-range case ($a=3,b=1.4$), the compute exponent improves monotonically with $\theta$ (Table~\ref{tab:compute}).  Weight-decay schedule experiments confirm Theorem~\ref{thm:adamw}: for $\delta(C)=C^{-s}$, the learned-mode exponent tracks $s/\alpha$ until the no-decay law is recovered.

\begin{table}[t]
\centering
\caption{Compute-optimal hard-source run with $a=3,b=1.4$.  The risk exponent improves monotonically with visible spectral information.}
\label{tab:compute}
\begin{tabular}{ccccc}
\toprule
$\theta$ & $\alpha=a(1-\theta/2)$ & observed risk exp. & predicted risk exp. & phase \\
\midrule
0.00 & 3.000 & $-0.102$ & $-0.100$ & time-limited \\
0.25 & 2.625 & $-0.113$ & $-0.110$ & time-limited \\
0.50 & 2.250 & $-0.128$ & $-0.123$ & time-limited \\
0.75 & 1.875 & $-0.148$ & $-0.139$ & time-limited \\
1.00 & 1.500 & $-0.167$ & $-0.160$ & time-limited \\
\bottomrule
\end{tabular}
\end{table}

\begin{figure}[t]
\centering
\fbox{\begin{minipage}{0.93\linewidth}
\textbf{Placeholder for Figure 3: compute-optimal phase diagram.}  Left: compute-risk exponent versus $\theta$, showing monotone improvement and saturation.  Middle: optimal $M,N,K$ exponents.  Right: AdamW schedule $s$ versus compute exponent, showing the cap-active and cap-inactive regimes.
\end{minipage}}
\caption{Visible spectral information changes compute-optimal scaling until the variance-limited phase; AdamW decay must shrink with compute to avoid a learned-mode cap.}
\label{fig:compute}
\end{figure}

\subsection{Random-feature and sketched-coordinate experiments}
Random-feature/sketch experiments test the bridge beyond diagonal coordinates.  In aligned and band-limited features, coordinatewise Adam/RMSProp/AdamW have $\qeff\approx0.48$--$0.49$.  In Haar and global Gaussian-sketch features, they have $\qeff\approx0$, while a spectral oracle remains at $\qeff=1/2$.

A wide learning-rate run resolves a key ambiguity: when the spectrum is hidden from coordinatewise second moments, a tuned spectral oracle achieves lower risk than coordinatewise Adam/RMSProp (Table~\ref{tab:randomfeatures}).  Thus global mixing does not remove the useful spectrum; it hides it from diagonal adaptivity.

\begin{table}[t]
\centering
\caption{Random-feature/sketch wide learning-rate run.  Global mixing hides spectral information from coordinatewise Adam/RMSProp, while the spectral oracle retains $\qeff=1/2$ and achieves lower tuned final risk.}
\label{tab:randomfeatures}
\begin{tabular}{llccc}
\toprule
case & optimizer & best LR & final risk & $\qeff$ \\
\midrule
Gaussian sketch & Adam & 0.01 & $8.35\times10^{-4}$ & $-0.028$ \\
Gaussian sketch & RMSProp & 0.01 & $1.52\times10^{-3}$ & $-0.032$ \\
Gaussian sketch & Spectral oracle & 0.10 & $1.45\times10^{-4}$ & $0.500$ \\
Haar & AdamW & 0.003 & $7.84\times10^{-4}$ & $-0.023$ \\
Haar & RMSProp & 0.003 & $9.50\times10^{-4}$ & $-0.027$ \\
Haar & Spectral oracle & 0.03 & $1.16\times10^{-4}$ & $0.500$ \\
\bottomrule
\end{tabular}
\end{table}

\section{Related work}

\paragraph{Scaling laws.}
Empirical scaling laws have been reported across language modeling, vision, generative modeling, and transfer \citep{hestness2017deep,rosenfeld2019constructive,kaplan2020scaling,henighan2020scaling,hernandez2021scaling,hoffmann2022training}.  Theoretical accounts connect power laws to kernel spectra, task-model alignment, and high-dimensional regression \citep{bordelon2020spectrum,canatar2021spectral,bahri2021explaining,lin2024scaling}.  Our closest baseline is \citet{lin2024scaling}, who prove model/data/compute scaling laws for one-pass SGD in sketched linear regression.  We study the complementary optimizer dimension: when preconditioning changes the spectrum, the scaling law itself changes.

\paragraph{Adaptive optimization and preconditioning.}
Adaptive and preconditioned methods include AdaGrad \citep{duchi2011adaptive}, RMSProp \citep{tieleman2012rmsprop}, Adam \citep{kingma2015adam}, AdamW \citep{loshchilov2019decoupled}, natural gradient \citep{amari1998natural}, K-FAC \citep{martens2015optimizing}, Shampoo \citep{gupta2018shampoo}, and related large-scale second-order methods \citep{anil2020scalable}.  Prior theory often focuses on convergence, regret, or implicit bias \citep{reddi2018convergence,wilson2017marginal}.  We instead ask whether the preconditioner reshapes a power-law spectrum enough to change scaling exponents.

\paragraph{Random features, kernels, and benign overfitting.}
Random features and kernel approximations provide tractable models of feature learning and high-dimensional generalization \citep{rahimi2007random,rudi2017generalization,bach2017breaking,jacot2018neural}.  Power-law eigenspectra also appear in studies of kernel learning curves and double descent \citep{belkin2019reconciling,hastie2022surprises,bartlett2020benign}.  Our feature-map results show that the same covariance spectrum can lead to different optimizer exponents depending on whether the trainable coordinate system exposes the spectrum to coordinatewise second moments.

\paragraph{Optimizer-dependent scaling.}
Recent empirical work has emphasized that optimizer choice and hyperparameter transfer can affect observed scaling behavior \citep{mccandlish2018empirical,smith2018bayesian,ramani2026optimizer}.  Our contribution is a minimal theorem-level mechanism: the visible spectrum controls the exponent.  This suggests reporting spectral diagnostics, not only final loss, when comparing optimizers across scale.

\section{Limitations and scope}

The theory is for Gaussian linear regression and feature-coordinate models.  Neural networks have non-Gaussian gradients, changing representations, layerwise normalization, and interactions between optimization and representation learning.  Our results should therefore be read as a mechanistic scaling-law model, not a complete theory of deep learning.  The experiments are synthetic by design and isolate one mechanism at a time.  Finally, finite-time risk improvements can occur even when $\qeff\approx0$; our claim concerns spectral scaling exponents, not all possible optimizer advantages.

\section{Conclusion}

Adaptive optimizers do not automatically change scaling laws.  They change scaling exponents when their coordinatewise second moments make the covariance spectrum visible.  The visible exponent $\theta$ determines $\qeff=\theta/2$, which determines the learned-mode count, risk law, compute-optimal allocation, and AdamW decay schedule.  This gives a concrete answer to when optimizer choice should enter scaling-law forecasts: not through adaptivity alone, but through the spectrum visible to the optimizer's preconditioner.

\appendix

\section{Proof details}

\subsection{Source stability}
\label{app:source}
If $P=f(\Sigma)$ is spectral, then $\tSig=P^{1/2}\Sigma P^{1/2}=\Sigma f(\Sigma)$ and $u^\star=P^{-1/2}w^\star$.  For each eigenvector $v_i$,
\[
    \ts_i
    =
    \lambda_i f(\lambda_i)\left(f(\lambda_i)^{-1/2}\langle v_i,w^\star\rangle\right)^2
    =s_i.
\]
For invariant band decompositions, total source mass is preserved blockwise:
\[
    (P_\ell^{-1/2}w^\star_\ell)^\top(P_\ell^{1/2}\Sigma_\ell P_\ell^{1/2})(P_\ell^{-1/2}w^\star_\ell)
    =
    (w^\star_\ell)^\top\Sigma_\ell w^\star_\ell.
\]
This is enough for band-averaged bias estimates because filters are constant up to factors inside comparable-eigenvalue bands.

\subsection{Risk filters}
The exact risk decomposition uses the bias and variance filters in \eqref{eq:filters}.  For two-slope spectra, with effective exponent $\alpha$ before the damping knee and exponent $a$ after it, integral comparison gives
\[
    \sum_{i\le M}\frac{\ts_i}{1+n\tmu_i}
    \asymp
    K(n)^{1-b}
\]
when $1<b<\alpha+1$ and $K(n)<M$.  If $\alpha>1/2$, then
\[
    \frac1{\Neff}\sum_{i\le M}\min\{1,n^2\tmu_i^2\}
    \asymp
    \frac{\min\{M,K(n)\}}{\Neff}.
\]
For Adam/RMSProp, $\alpha=a/2$, so $\alpha>1/2$ follows from $a>1$.

\subsection{Compute-optimal derivation}
For fixed $N$, minimize $K^{1-b}+K/N$ under $K\le N^{1/\alpha}$.  The unconstrained optimum satisfies $K^{-b}\asymp N^{-1}$, hence $K\asymp N^{1/b}$.  If $b\le\alpha$, this exceeds the time cap, so $K\asymp N^{1/\alpha}$; otherwise $K\asymp N^{1/b}$.  Combining with $M=C/N$ and balancing $M^{1-b}$ against the optimized $N$ term yields Theorem~\ref{thm:compute}.

\subsection{Diagnostic consistency}
Suppose on a fitting window $i\in[L,U]$,
\[
    \lambda_i=c_\lambda i^{-a}e^{\varepsilon_i},
    \qquad
    \mu_i=c_\mu i^{-\alpha}e^{\delta_i},
\]
with $|\varepsilon_i|,|\delta_i|\le\eta$.  Least-squares log-log slope estimates satisfy
\[
    |\widehat a-a|+|\widehat\alpha-\alpha|
    \lesssim
    \eta/\log(U/L).
\]
Thus $\widehat q=1-\widehat\alpha/\widehat a$ is consistent whenever $a$ is bounded away from zero and the fitting window has nontrivial logarithmic width.  This justifies reporting $\qeff$ as the primary mechanism diagnostic.

\subsection{Hidden-spectrum oracle separation}
Suppose coordinatewise adaptivity has $q_{\rm coord}=0$ while a spectral oracle has $q_{\rm spec}=1/2$.  Let
\[
    \beta(q)=\frac{b-1}{\max\{a(1-q),b\}+1}.
\]
If $\beta(1/2)>\beta(0)$, then
\[
    \frac{R^\star_{q=0}(C)-\sigma^2}{R^\star_{q=1/2}(C)-\sigma^2}
    \asymp
    C^{\beta(1/2)-\beta(0)}\to\infty.
\]
Scalar learning-rate tuning rescales the effective horizon by a constant and cannot change $\beta(q)$.  This explains why coordinatewise Adam/RMSProp can improve constants in Haar or Gaussian-sketch coordinates while remaining in the $q=0$ exponent class.

\section{Additional experiment details}

\paragraph{Deterministic filters.}  We evaluate the spectral filters $B(n,M)$ and $V(n,M)$ directly for power-law spectra and fit slopes over geometric ranges of $n$.

\paragraph{Stochastic training.}  We train mini-batch Gaussian linear regression with SGD, RMSProp, Adam, AdamW, coordinate oracles, and spectral oracles.  For adaptive methods, $\qeff$ is estimated by the slope change between the covariance spectrum and the effective covariance induced by the current second-moment preconditioner.

\paragraph{Coordinate and feature cases.}  Aligned uses the covariance eigenbasis.  Band-limited features apply random rotations within comparable-eigenvalue blocks.  Flat/Haar/global Gaussian-sketch features mix spectral scales broadly and make coordinate variances nearly scalar.

\paragraph{Compute sweeps.}  Compute-optimal experiments optimize the deterministic proxy \eqref{eq:compute-proxy} over $(M,N,K)$ for each compute budget.  AdamW schedule experiments add the cap $K\le\delta(C)^{-1/\alpha}$.

\begin{thebibliography}{99}

\bibitem[Amari(1998)]{amari1998natural}
Shun-ichi Amari.
\newblock Natural gradient works efficiently in learning.
\newblock \emph{Neural Computation}, 10(2):251--276, 1998.

\bibitem[Anil et~al.(2020)]{anil2020scalable}
Rohan Anil, Vineet Gupta, Tomer Koren, and Yoram Singer.
\newblock Scalable second order optimization for deep learning.
\newblock \emph{arXiv preprint arXiv:2002.09018}, 2020.

\bibitem[Bach(2017)]{bach2017breaking}
Francis Bach.
\newblock Breaking the curse of dimensionality with convex neural networks.
\newblock \emph{Journal of Machine Learning Research}, 18(19):1--53, 2017.

\bibitem[Bahri et~al.(2021)]{bahri2021explaining}
Yasaman Bahri, Ethan Dyer, Jared Kaplan, Jaehoon Lee, and Utkarsh Sharma.
\newblock Explaining neural scaling laws.
\newblock \emph{arXiv preprint arXiv:2102.06701}, 2021.

\bibitem[Bartlett et~al.(2020)]{bartlett2020benign}
Peter~L. Bartlett, Philip~M. Long, Gabor Lugosi, and Alexander Tsigler.
\newblock Benign overfitting in linear regression.
\newblock \emph{Proceedings of the National Academy of Sciences}, 117(48):30063--30070, 2020.

\bibitem[Belkin et~al.(2019)]{belkin2019reconciling}
Mikhail Belkin, Daniel Hsu, Siyuan Ma, and Soumik Mandal.
\newblock Reconciling modern machine-learning practice and the classical bias--variance trade-off.
\newblock \emph{Proceedings of the National Academy of Sciences}, 116(32):15849--15854, 2019.

\bibitem[Bordelon et~al.(2020)]{bordelon2020spectrum}
Blake Bordelon, Abdulkadir Canatar, and Cengiz Pehlevan.
\newblock Spectrum dependent learning curves in kernel regression and wide neural networks.
\newblock \emph{International Conference on Machine Learning}, 2020.

\bibitem[Canatar et~al.(2021)]{canatar2021spectral}
Abdulkadir Canatar, Blake Bordelon, and Cengiz Pehlevan.
\newblock Spectral bias and task-model alignment explain generalization in kernel regression and infinitely wide neural networks.
\newblock \emph{Nature Communications}, 12:2914, 2021.

\bibitem[Duchi et~al.(2011)]{duchi2011adaptive}
John Duchi, Elad Hazan, and Yoram Singer.
\newblock Adaptive subgradient methods for online learning and stochastic optimization.
\newblock \emph{Journal of Machine Learning Research}, 12:2121--2159, 2011.

\bibitem[Gupta et~al.(2018)]{gupta2018shampoo}
Vineet Gupta, Tomer Koren, and Yoram Singer.
\newblock Shampoo: Preconditioned stochastic tensor optimization.
\newblock \emph{International Conference on Machine Learning}, 2018.

\bibitem[Hastie et~al.(2022)]{hastie2022surprises}
Trevor Hastie, Andrea Montanari, Saharon Rosset, and Ryan~J. Tibshirani.
\newblock Surprises in high-dimensional ridgeless least squares interpolation.
\newblock \emph{Annals of Statistics}, 50(2):949--986, 2022.

\bibitem[Henighan et~al.(2020)]{henighan2020scaling}
Tom Henighan, Jared Kaplan, Mor Katz, Mark Chen, Christopher Hesse, Jacob Jackson, Heewoo Jun, Tom~B. Brown, Prafulla Dhariwal, Scott Gray, Chris Hallacy, Benjamin Mann, Alec Radford, Aditya Ramesh, Nick Ryder, Daniel~M. Ziegler, John Schulman, Dario Amodei, and Sam McCandlish.
\newblock Scaling laws for autoregressive generative modeling.
\newblock \emph{arXiv preprint arXiv:2010.14701}, 2020.

\bibitem[Hernandez et~al.(2021)]{hernandez2021scaling}
Danny Hernandez, Jared Kaplan, Tom Henighan, and Sam McCandlish.
\newblock Scaling laws for transfer.
\newblock \emph{arXiv preprint arXiv:2102.01293}, 2021.

\bibitem[Hestness et~al.(2017)]{hestness2017deep}
Joel Hestness, Sharan Narang, Newsha Ardalani, Gregory Diamos, Heewoo Jun, Hassan Kianinejad, Md~Mostofa Ali Patwary, Yang Yang, and Yanqi Zhou.
\newblock Deep learning scaling is predictable, empirically.
\newblock \emph{arXiv preprint arXiv:1712.00409}, 2017.

\bibitem[Hoffmann et~al.(2022)]{hoffmann2022training}
Jordan Hoffmann, Sebastian Borgeaud, Arthur Mensch, Elena Buchatskaya, Trevor Cai, Eliza Rutherford, Diego de~Las Casas, Lisa~Anne Hendricks, Johannes Welbl, Aidan Clark, Tom Hennigan, Eric Noland, Katie Millican, George van~den Driessche, Bogdan Damoc, Aurelia Guy, Simon Osindero, Karen Simonyan, Erich Elsen, Jack~W. Rae, Oriol Vinyals, and Laurent Sifre.
\newblock Training compute-optimal large language models.
\newblock \emph{arXiv preprint arXiv:2203.15556}, 2022.

\bibitem[Jacot et~al.(2018)]{jacot2018neural}
Arthur Jacot, Franck Gabriel, and Clement Hongler.
\newblock Neural tangent kernel: Convergence and generalization in neural networks.
\newblock \emph{Advances in Neural Information Processing Systems}, 2018.

\bibitem[Kaplan et~al.(2020)]{kaplan2020scaling}
Jared Kaplan, Sam McCandlish, Tom Henighan, Tom~B. Brown, Benjamin Chess, Rewon Child, Scott Gray, Alec Radford, Jeffrey Wu, and Dario Amodei.
\newblock Scaling laws for neural language models.
\newblock \emph{arXiv preprint arXiv:2001.08361}, 2020.

\bibitem[Kingma and Ba(2015)]{kingma2015adam}
Diederik~P. Kingma and Jimmy Ba.
\newblock Adam: A method for stochastic optimization.
\newblock \emph{International Conference on Learning Representations}, 2015.

\bibitem[Lin et~al.(2024)]{lin2024scaling}
Licong Lin, Jingfeng Wu, Sham~M. Kakade, Peter~L. Bartlett, and Jason~D. Lee.
\newblock Scaling laws in linear regression: Compute, parameters, and data.
\newblock \emph{Advances in Neural Information Processing Systems}, 2024.

\bibitem[Lin et~al.(2025)]{lin2025datareuse}
Licong Lin, Jingfeng Wu, and Peter~L. Bartlett.
\newblock Improved scaling laws in linear regression via data reuse.
\newblock \emph{arXiv preprint arXiv:2506.08415}, 2025.

\bibitem[Loshchilov and Hutter(2019)]{loshchilov2019decoupled}
Ilya Loshchilov and Frank Hutter.
\newblock Decoupled weight decay regularization.
\newblock \emph{International Conference on Learning Representations}, 2019.

\bibitem[Martens and Grosse(2015)]{martens2015optimizing}
James Martens and Roger Grosse.
\newblock Optimizing neural networks with Kronecker-factored approximate curvature.
\newblock \emph{International Conference on Machine Learning}, 2015.

\bibitem[McCandlish et~al.(2018)]{mccandlish2018empirical}
Sam McCandlish, Jared Kaplan, Dario Amodei, and OpenAI Dota Team.
\newblock An empirical model of large-batch training.
\newblock \emph{arXiv preprint arXiv:1812.06162}, 2018.

\bibitem[Polyak(1964)]{polyak1964some}
Boris~T. Polyak.
\newblock Some methods of speeding up the convergence of iteration methods.
\newblock \emph{USSR Computational Mathematics and Mathematical Physics}, 4(5):1--17, 1964.

\bibitem[Rahimi and Recht(2007)]{rahimi2007random}
Ali Rahimi and Benjamin Recht.
\newblock Random features for large-scale kernel machines.
\newblock \emph{Advances in Neural Information Processing Systems}, 2007.

\bibitem[Ramani and Jain(2026)]{ramani2026optimizer}
Vansh Ramani and Shourya Vir Jain.
\newblock On the optimizer dependence of neural scaling laws.
\newblock \emph{Preprint}, 2026.

\bibitem[Reddi et~al.(2018)]{reddi2018convergence}
Sashank~J. Reddi, Satyen Kale, and Sanjiv Kumar.
\newblock On the convergence of Adam and beyond.
\newblock \emph{International Conference on Learning Representations}, 2018.

\bibitem[Rosenfeld et~al.(2019)]{rosenfeld2019constructive}
Jonathan~S. Rosenfeld, Amir Rosenfeld, Yonatan Belinkov, and Nir Shavit.
\newblock A constructive prediction of the generalization error across scales.
\newblock \emph{International Conference on Learning Representations}, 2020.

\bibitem[Rudi and Rosasco(2017)]{rudi2017generalization}
Alessandro Rudi and Lorenzo Rosasco.
\newblock Generalization properties of learning with random features.
\newblock \emph{Advances in Neural Information Processing Systems}, 2017.

\bibitem[Smith and Le(2018)]{smith2018bayesian}
Samuel~L. Smith and Quoc~V. Le.
\newblock A Bayesian perspective on generalization and stochastic gradient descent.
\newblock \emph{International Conference on Learning Representations}, 2018.

\bibitem[Tieleman and Hinton(2012)]{tieleman2012rmsprop}
Tijmen Tieleman and Geoffrey Hinton.
\newblock Lecture 6.5---RMSProp: Divide the gradient by a running average of its recent magnitude.
\newblock \emph{COURSERA: Neural Networks for Machine Learning}, 2012.

\bibitem[Wilson et~al.(2017)]{wilson2017marginal}
Ashia~C. Wilson, Rebecca Roelofs, Mitchell Stern, Nathan Srebro, and Benjamin Recht.
\newblock The marginal value of adaptive gradient methods in machine learning.
\newblock \emph{Advances in Neural Information Processing Systems}, 2017.

\end{thebibliography}

\end{document}
